% Options for packages loaded elsewhere
\PassOptionsToPackage{unicode}{hyperref}
\PassOptionsToPackage{hyphens}{url}
\documentclass[
]{article}
\usepackage{xcolor}
\usepackage[margin=1in]{geometry}
\usepackage{amsmath,amssymb}
\setcounter{secnumdepth}{-\maxdimen} % remove section numbering
\usepackage{iftex}
\ifPDFTeX
  \usepackage[T1]{fontenc}
  \usepackage[utf8]{inputenc}
  \usepackage{textcomp} % provide euro and other symbols
\else % if luatex or xetex
  \usepackage{unicode-math} % this also loads fontspec
  \defaultfontfeatures{Scale=MatchLowercase}
  \defaultfontfeatures[\rmfamily]{Ligatures=TeX,Scale=1}
\fi
\usepackage{lmodern}
\ifPDFTeX\else
  % xetex/luatex font selection
\fi
% Use upquote if available, for straight quotes in verbatim environments
\IfFileExists{upquote.sty}{\usepackage{upquote}}{}
\IfFileExists{microtype.sty}{% use microtype if available
  \usepackage[]{microtype}
  \UseMicrotypeSet[protrusion]{basicmath} % disable protrusion for tt fonts
}{}
\makeatletter
\@ifundefined{KOMAClassName}{% if non-KOMA class
  \IfFileExists{parskip.sty}{%
    \usepackage{parskip}
  }{% else
    \setlength{\parindent}{0pt}
    \setlength{\parskip}{6pt plus 2pt minus 1pt}}
}{% if KOMA class
  \KOMAoptions{parskip=half}}
\makeatother
\usepackage{graphicx}
\makeatletter
\newsavebox\pandoc@box
\newcommand*\pandocbounded[1]{% scales image to fit in text height/width
  \sbox\pandoc@box{#1}%
  \Gscale@div\@tempa{\textheight}{\dimexpr\ht\pandoc@box+\dp\pandoc@box\relax}%
  \Gscale@div\@tempb{\linewidth}{\wd\pandoc@box}%
  \ifdim\@tempb\p@<\@tempa\p@\let\@tempa\@tempb\fi% select the smaller of both
  \ifdim\@tempa\p@<\p@\scalebox{\@tempa}{\usebox\pandoc@box}%
  \else\usebox{\pandoc@box}%
  \fi%
}
% Set default figure placement to htbp
\def\fps@figure{htbp}
\makeatother
\setlength{\emergencystretch}{3em} % prevent overfull lines
\providecommand{\tightlist}{%
  \setlength{\itemsep}{0pt}\setlength{\parskip}{0pt}}
\usepackage{bookmark}
\IfFileExists{xurl.sty}{\usepackage{xurl}}{} % add URL line breaks if available
\urlstyle{same}
\hypersetup{
  pdftitle={Séries Temporais para Ciências Atuariais},
  pdfauthor={Marcos R. Gonzaga},
  hidelinks,
  pdfcreator={LaTeX via pandoc}}

\title{Séries Temporais para Ciências Atuariais}
\author{Marcos R. Gonzaga}
\date{2026-03-13}

\begin{document}
\maketitle

\subsection{Introdução}\label{introduuxe7uxe3o}

Este material apresenta conceitos, definições e exemplos práticos sobre
a metodologia de séries temporaos applicada às ciências atuariais..

\subsection{Séries Temporais
(Definição)}\label{suxe9ries-temporais-definiuxe7uxe3o}

Entender Séries Temporais é como aprender a ler o ``ritmo'' dos dados.
Diferente de uma análise convencional onde os dados são independentes,
aqui a ordem dos fatores altera (e muito) o produto.

\textbf{O Conceito: O que diferencia uma Série Temporal?}

Uma série temporal é um conjunto de observações coletadas em intervalos
de tempo sucessivos e (geralmente) iguais. O segredo está na
\emph{autocorrelação}: o valor de hoje depende, em algum nível, do valor
de ontem.

Todo modelo de séries temporais é divido em quatro componentes
principais:

\textbf{Tendência}: A direção de longo prazo (está subindo ou
descendo?).

\textbf{Sazonalidade}: Padrões que se repetem em intervalos fixos (ex:
vendas de sorvete no verão).

\textbf{Ciclo}: Flutuações de longo prazo sem período fixo (ex: ciclos
econômicos).

\textbf{Ruído (Erro)}: A variação aleatória que não conseguimos
explicar.

\textbf{O modelo na forma geral:}

\[Z_t=T+S+E\hspace{3cm}(1)\]

onde:

\begin{itemize}
\item
  \(Z_t\) é o valor da série em cada ponto do tempo t;
\item
  \(T\) é a tendência (determinística/estocástica) que a série possui ao
  longo do tempo;
\item
  \(E\) é o resíduo ou o erro aleatório do modelo.
\end{itemize}

\section{Para começar}\label{para-comeuxe7ar}

\textbf{1. Carregue os dados:}

data(``AirPassengers'') plot(AirPassengers, main=``Passageiros de Linhas
Aéreas (1949-1960)'', ylab=``Total de Passageiros'', xlab=``Ano'',
col=``blue'')

\end{document}
